\chapter{募资从零到一：材料、路演与合规边界}

\section{材料体系：PPM、数据室与 DDQ}

私募配售备忘录（PPM / 私募备忘录）是面向合格投资者的\textbf{核心法律与营销混合文件}，需披露策略、风险因素、费用、利益冲突、管理团队与历史业绩（若有）。\textbf{数据室（VDR）}存放基金章程、LPA 草案、合规政策样本、过往审计与参考投资组合（在允许范围内）。\textbf{尽职调查问卷（DDQ）}回复机构 LP 的标准问题，便于进入 allocator 的内部评分流程。三类材料应\textbf{同源更新}，避免路演口述与书面文件不一致引发监管或民事风险。

\section{路演节奏与叙事}

首轮沟通通常以\textbf{一页 Teaser}与简短电话开始，确认策略与规模是否匹配 allocator 的授权。正式路演聚焦：投资逻辑可重复性、团队分工、风险控制、费用与条款、与同类基金的差异化。\textbf{演示材料}须经合规审查（若管理人已持牌或在持牌顾问指导下）。问答环节应如实记录；对无法承诺的事项（如未来收益）明确拒绝书面化。

\section{合格投资者与宣传边界}

各法域对「合格投资者」「专业客户」的定义不同；本书合规篇提供框架性索引。一般原则是：\textbf{非公开募集}、\textbf{不对不特定对象宣传保本或确定收益}、\textbf{适当性匹配}。社交媒体与公开会议上的表述极易构成「公开劝诱」的争议，冷启动章节亦强调有节制的触达；本章从基金文件角度要求：所有对外版本留有\textbf{版本号与审批记录}。

\section{首次交割与后续关闭}

首次关闭（first closing）确定首批 LP 的认缴与计费起点；后续关闭（subsequent closings）可能涉及\textbf{同等条款最惠}、对早期 LP 的补偿（equalisation）与对价调整。行政上需协调认购书签署、KYC、资金到账与工商/登记变更（视载体而定）。延长的募集期会增加法律与市场成本，应在 PPM 中设定明确截止条件与 GP 的权利。
